\documentclass[11pt]{ctexart}
\usepackage[a4paper,margin=1in]{geometry}
\usepackage{longtable,booktabs}
\title{Geant4-Agent Casebank Regression Report (EN)}
\author{Automated Pipeline}
\date{2026-03-07\_172117}
\begin{document}
\maketitle
\section{Unified Pass/Fail Table}
\begin{longtable}{p{0.06\linewidth}p{0.08\linewidth}p{0.07\linewidth}p{0.08\linewidth}p{0.07\linewidth}p{0.07\linewidth}p{0.07\linewidth}p{0.07\linewidth}p{0.09\linewidth}p{0.06\linewidth}p{0.21\linewidth}}
\toprule
ID & Domain & Style & Completeness & ExpInit & ActInit & ExpFinal & ActFinal & Missing(I->F) & Result & Fail Reason \\
\midrule
MTM01 & medical & fuzzy & missing & FAIL & FAIL & PASS & PASS & 12->0 & PASS & - \\
MTM02 & aerospace & precise & missing & FAIL & FAIL & PASS & PASS & 2->0 & PASS & - \\
MTM03 & industrial\_ndt & precise & missing & FAIL & FAIL & PASS & PASS & 12->0 & PASS & - \\
MTM04 & physics & precise & missing & FAIL & FAIL & PASS & PASS & 8->0 & PASS & - \\
MTM05 & nuclear\_engineering & fuzzy & missing & FAIL & FAIL & PASS & PASS & 11->0 & PASS & - \\
MTM06 & security\_screening & fuzzy & missing & FAIL & FAIL & PASS & PASS & 11->0 & PASS & - \\
MTM07 & semiconductor & precise & missing & FAIL & FAIL & PASS & PASS & 3->0 & PASS & - \\
MTM08 & education & fuzzy & missing & FAIL & FAIL & PASS & PASS & 11->0 & PASS & - \\
MTM09 & environmental\_radiation & fuzzy & missing & FAIL & FAIL & PASS & PASS & 3->0 & PASS & - \\
MTM10 & high\_energy\_physics & precise & missing & FAIL & FAIL & PASS & PASS & 10->0 & PASS & - \\
MTM11 & materials\_update & precise & full & PASS & PASS & PASS & PASS & 0->0 & PASS & - \\
MTM12 & overwrite\_reject & precise & full & PASS & PASS & PASS & PASS & 0->0 & PASS & - \\
MTM13 & medical\_detector & fuzzy & missing & FAIL & FAIL & PASS & PASS & 11->0 & FAIL & expected\_final\_mismatch:geometry.structure \\
MTM14 & industrial\_panel & precise & missing & FAIL & FAIL & PASS & PASS & 16->0 & FAIL & expected\_final\_mismatch:geometry.structure \\
MTM15 & shielding\_stack & precise & missing & FAIL & FAIL & PASS & FAIL & 15->6 & FAIL & final\_complete\_mismatch, expected\_final\_mismatch:geometry.structure \\
MTM16 & nested\_target & precise & missing & FAIL & FAIL & PASS & FAIL & 11->9 & FAIL & final\_complete\_mismatch, expected\_final\_mismatch:geometry.structure, expected\_final\_mismatch:materials.selected\_materials.0, expected\_final\_mismatch:source.type, expected\_final\_mismatch:source.particle, expected\_final\_mismatch:output.format \\
MTM17 & coaxial\_shield & fuzzy & missing & FAIL & FAIL & PASS & FAIL & 9->3 & FAIL & final\_complete\_mismatch, expected\_final\_mismatch:geometry.structure \\
MTM18 & boolean\_cavity & precise & missing & FAIL & FAIL & PASS & PASS & 12->0 & FAIL & expected\_final\_mismatch:geometry.structure \\
MTM19 & ring\_complete & precise & full & PASS & FAIL & PASS & FAIL & 2->2 & FAIL & initial\_complete\_mismatch, final\_complete\_mismatch \\
MTM20 & grid\_complete & precise & full & PASS & PASS & PASS & PASS & 0->0 & FAIL & expected\_final\_mismatch:materials.selected\_materials.0 \\
MTM21 & layered\_detector & fuzzy & missing & FAIL & FAIL & PASS & PASS & 15->0 & FAIL & expected\_final\_mismatch:geometry.structure \\
MTM22 & embedded\_sensor & precise & full & PASS & FAIL & PASS & FAIL & 13->13 & FAIL & initial\_complete\_mismatch, final\_complete\_mismatch, expected\_final\_mismatch:geometry.structure, expected\_final\_mismatch:materials.selected\_materials.0, expected\_final\_mismatch:source.type, expected\_final\_mismatch:output.format \\
MTM23 & layered\_shield\_complete & precise & full & PASS & PASS & PASS & PASS & 0->0 & FAIL & expected\_final\_mismatch:geometry.structure \\
MTM24 & boolean\_complete & precise & full & PASS & PASS & PASS & PASS & 0->0 & PASS & - \\
\bottomrule
\end{longtable}
\section{Baseline Comparison}
\begin{itemize}
\item pass\_rate=0.542
\item initial\_alignment\_rate=0.917
\item final\_alignment\_rate=0.792
\item avg\_turns\_per\_case=2.375
\item missing\_reduction\_rate=1.000
\item confirm\_apply\_success\_rate=1.000
\item rollback\_after\_confirm\_rate=0.000
\item internal\_leak\_turn\_rate=0.123
\item repeated\_assistant\_case\_rate=0.000
\item stale\_confirm\_case\_rate=0.000
\item delta\_pass\_rate=-0.408
\item regressions=none
\item improvements=none
\end{itemize}
\appendix
\section{Appendix: Full Dialogues and Naturalized Internal Chains}
\subsection{Case MTM01 (medical)}
\textbf{Dialogue:}
\begin{itemize}
\item User: 我先要一个用于医学成像的初步辐照场景，细节还没定，你先告诉我还缺什么。
\item Assistant: 好的，我们目前已经确定了使用“盒体”作为几何结构。

接下来还需要您提供以下信息才能完成配置：
1.  盒体在x、y、z三个方向的具体尺寸。
2.  这个体积需要绑定什么材料。
3.  您希望使用什么类型的放射源（例如伽马源、电子源等）。
\item User: 采用 300 mm x 300 mm x 300 mm 的水箱，材料 G4\_WATER；源用束流 gamma，能量 6 MeV，位置 (0,0,-400) mm，方向 +z；物理列表 QBBC；输出 root。
\item Assistant: 我注意到这可能会覆盖已经确认的设置。具体来说：几何结构类型将从“未知”改为“单盒体”；盒体x方向尺寸将设置为300.0。

如果您确认应用这些修改，请直接回复“确认”；如果您希望保留原来的值，请回复“保持原值”。
\item User: 确认
\item Assistant: 好的，配置已完成！目前已经确认的设置包括：几何结构类型，以及盒体在x、y、z三个方向上的尺寸。
\end{itemize}
\textbf{Internal Chain (Natural Language):}
\begin{itemize}
\item Turn 1: action=summarize\_progress, phase=geometry, complete=False.
\item Missing fields: geometry.params.module\_x, geometry.params.module\_y, geometry.params.module\_z, materials.volume\_material\_map, source.type, source.particle, source.energy, source.position, source.direction, physics.physics\_list, output.format, output.path.
\item LLM used=True; fallback=none; backend=llm\_semantic\_frame+runtime\_semantic.
\item Turn 2: action=confirm\_overwrite, phase=geometry, complete=False.
\item Missing fields: geometry.params.module\_x, geometry.params.module\_y, geometry.params.module\_z.
\item LLM used=True; fallback=none; backend=llm\_slot\_frame+runtime\_semantic.
\item Turn 3: action=finalize, phase=materials, complete=True.
\item Missing fields: none.
\item LLM used=True; fallback=none; backend=llm\_slot\_frame+runtime\_semantic+confirmed\_pending\_overwrite.
\end{itemize}
\subsection{Case MTM02 (aerospace)}
\textbf{Dialogue:}
\begin{itemize}
\item User: I need a first-pass proton irradiation setup for a silicon device; keep it simple for now.
\item Assistant: Geometry: updated geometry structure, box size along x, box size along y. Materials: updated selected materials, volume-to-material binding. Source: updated source type, particle type, source energy. Physics: updated physics list. Output: still needs output format, output path.
\item User: Use a silicon box 50 mm x 50 mm x 5 mm, beam source proton 100 MeV from (0,0,-100) mm along +z, physics FTFP\_BERT, output json.
\item Assistant: You've requested a change that would overwrite a confirmed value.

Pending changes:
- Box size along x: 1000.0 mm → 50.0 mm
- Box size along y: 1000.0 mm → 50.0 mm

Please reply **"confirm"** to apply these changes, or **"keep existing"** to retain the current values.
\item User: confirm
\item Assistant: Your configuration is complete. The following items are confirmed: source type, box size along z, box size along y, and box size along x.
\end{itemize}
\textbf{Internal Chain (Natural Language):}
\begin{itemize}
\item Turn 1: action=summarize\_progress, phase=materials, complete=False.
\item Missing fields: output.format, output.path.
\item LLM used=True; fallback=none; backend=llm\_slot\_frame+runtime\_semantic.
\item Turn 2: action=confirm\_overwrite, phase=source, complete=False.
\item Missing fields: none.
\item LLM used=True; fallback=none; backend=llm\_slot\_frame+runtime\_semantic.
\item Turn 3: action=finalize, phase=physics, complete=True.
\item Missing fields: none.
\item LLM used=True; fallback=none; backend=llm\_semantic\_frame+runtime\_semantic+confirmed\_pending\_overwrite.
\end{itemize}
\subsection{Case MTM03 (industrial\_ndt)}
\textbf{Dialogue:}
\begin{itemize}
\item User: 我要做一个 gamma 穿透铁靶的小测试，但我还没把完整参数说完。
\item Assistant: 目前我们已确认了材料和粒子类型（gamma 射线与铁靶）。接下来还需要补充几何部分的参数（结构类型、盒体尺寸），以及体积与材料绑定的设置。请继续告诉我这些信息，我们就能继续推进配置了。
\item User: 请设为半径 40 mm、半长 80 mm 的铁圆柱，材料 G4\_Fe；点源 gamma，能量 2 MeV，位置 (0,0,-200) mm，方向 +z；物理列表 Shielding；输出 csv。
\item Assistant: 配置已完成。当前已确认的参数包括：体积与材料绑定、输出路径、输出格式和物理列表。
\end{itemize}
\textbf{Internal Chain (Natural Language):}
\begin{itemize}
\item Turn 1: action=summarize\_progress, phase=geometry, complete=False.
\item Missing fields: geometry.structure, geometry.params.module\_x, geometry.params.module\_y, geometry.params.module\_z, materials.volume\_material\_map, source.type, source.energy, source.position, source.direction, physics.physics\_list, output.format, output.path.
\item LLM used=True; fallback=none; backend=llm\_slot\_frame+runtime\_semantic.
\item Turn 2: action=finalize, phase=materials, complete=True.
\item Missing fields: none.
\item LLM used=True; fallback=none; backend=llm\_slot\_frame+runtime\_semantic.
\end{itemize}
\subsection{Case MTM04 (physics)}
\textbf{Dialogue:}
\begin{itemize}
\item User: Gamma point source 1 MeV at (0,0,-100) toward +z; please build the rest.
\item Assistant: Geometry: still needs geometry structure, box size along x. Materials: still needs volume-to-material binding. Source: updated source type, particle type, source energy. Physics: still needs physics list. Output: still needs output format, output path.
\item User: Set copper box 1 m x 1 m x 1 m, physics FTFP\_BERT, output hdf5.
\item Assistant: Your configuration is complete. The following items are confirmed: volume-to-material binding, output path, output format, and physics list.
\end{itemize}
\textbf{Internal Chain (Natural Language):}
\begin{itemize}
\item Turn 1: action=summarize\_progress, phase=geometry, complete=False.
\item Missing fields: geometry.structure, geometry.params.module\_x, geometry.params.module\_y, geometry.params.module\_z, materials.volume\_material\_map, physics.physics\_list, output.format, output.path.
\item LLM used=True; fallback=none; backend=llm\_slot\_frame+runtime\_semantic.
\item Turn 2: action=finalize, phase=materials, complete=True.
\item Missing fields: none.
\item LLM used=True; fallback=none; backend=llm\_slot\_frame+runtime\_semantic.
\end{itemize}
\subsection{Case MTM05 (nuclear\_engineering)}
\textbf{Dialogue:}
\begin{itemize}
\item User: 我想做一个中子屏蔽的概念验证，现在先给一个缺参开头。
\item Assistant: 好的，我们已确认了“几何结构类型”和“粒子类型”。接下来还需要您补充以下信息：

**几何**：盒体在x、y、z三个方向的具体尺寸。
**材料**：体积与材料的绑定关系。
**源**：源的类型（例如，点源、面源等）。

请随时告诉我这些细节，我们继续配置。
\item User: 使用 200 mm x 200 mm x 200 mm 的铅块，材料 G4\_Pb；点源 neutron，能量 14 MeV，位置 (0,0,-50) mm，方向 +z；物理列表 QGSP\_BERT\_HP；输出 xml。
\item Assistant: 我检测到这会覆盖已经确认的值。待修改内容：几何结构类型：unknown -> single\_box；盒体 x 方向尺寸：None -> 200.0。如果要应用这些修改，请直接回复“确认”；如果要保留原值，请直接说“保持原值”。
\item User: 确认
\item Assistant: 配置已完成！目前盒体的尺寸（x、y、z方向）和几何结构类型都已确认。
\end{itemize}
\textbf{Internal Chain (Natural Language):}
\begin{itemize}
\item Turn 1: action=summarize\_progress, phase=geometry, complete=False.
\item Missing fields: geometry.params.module\_x, geometry.params.module\_y, geometry.params.module\_z, materials.volume\_material\_map, source.type, source.energy, source.position, source.direction, physics.physics\_list, output.format, output.path.
\item LLM used=True; fallback=none; backend=llm\_slot\_frame+runtime\_semantic.
\item Turn 2: action=confirm\_overwrite, phase=geometry, complete=False.
\item Missing fields: geometry.params.module\_x, geometry.params.module\_y, geometry.params.module\_z.
\item LLM used=True; fallback=none; backend=llm\_slot\_frame+runtime\_semantic.
\item Turn 3: action=finalize, phase=materials, complete=True.
\item Missing fields: none.
\item LLM used=True; fallback=none; backend=llm\_semantic\_frame+runtime\_semantic+confirmed\_pending\_overwrite.
\end{itemize}
\subsection{Case MTM06 (security\_screening)}
\textbf{Dialogue:}
\begin{itemize}
\item User: Need a quick cargo-screening toy setup with forward photons, but I have not fixed the target shape, material, source energy, or output format yet.
\item Assistant: Geometry: updated geometry structure; still needs box size along x, box size along y. Materials: still needs volume-to-material binding. Source: updated particle type; still needs source type, source energy. Physics: still needs physics list. Output: still needs output format, output path.
\item User: Use steel cylinder target radius 60 mm half-length 120 mm, beam gamma 3 MeV from (0,0,-500) mm to +z, physics Shielding, output root.
\item Assistant: This would overwrite an already confirmed value. Pending change: geometry structure: unknown -> single\_tubs; radius: None -> 60.0. Reply 'confirm' to apply it, or say 'keep existing' to keep the current value.
\item User: confirm
\item Assistant: Configuration complete. The following items are confirmed: box size along z, box size along y, box size along x, and volume-to-material binding.
\end{itemize}
\textbf{Internal Chain (Natural Language):}
\begin{itemize}
\item Turn 1: action=summarize\_progress, phase=geometry, complete=False.
\item Missing fields: geometry.params.module\_x, geometry.params.module\_y, geometry.params.module\_z, materials.volume\_material\_map, source.type, source.energy, source.position, source.direction, physics.physics\_list, output.format, output.path.
\item LLM used=True; fallback=none; backend=llm\_slot\_frame+runtime\_semantic.
\item Turn 2: action=confirm\_overwrite, phase=geometry, complete=False.
\item Missing fields: geometry.params.module\_x, geometry.params.module\_y, geometry.params.module\_z.
\item LLM used=True; fallback=none; backend=llm\_slot\_frame+runtime\_semantic.
\item Turn 3: action=finalize, phase=materials, complete=True.
\item Missing fields: none.
\item LLM used=True; fallback=none; backend=llm\_semantic\_frame+runtime\_semantic+confirmed\_pending\_overwrite.
\end{itemize}
\subsection{Case MTM07 (semiconductor)}
\textbf{Dialogue:}
\begin{itemize}
\item User: 我想搭一个硅薄片的中子照射测试，但先只给一个简述。
\item Assistant: 好的，我们已同步了以下配置：
*   **几何**：盒体结构及其尺寸。
*   **材料**：已选定材料并完成绑定。
*   **源**：粒子类型、能量和位置。

接下来，还需要您补充以下信息才能完成配置：
*   **物理列表**
*   **输出格式**
*   **输出路径**
\item User: 使用 20 mm x 20 mm x 1 mm 的硅片，材料 G4\_Si；各向同性 neutron，能量 2 MeV，位置 (0,0,0) mm；物理列表 QGSP\_BERT\_HP；输出 json。
\item Assistant: 我注意到这将会覆盖已经确认的配置项。具体修改如下：

*   **盒体 x 方向尺寸**：从 1000.0 改为 20.0
*   **盒体 y 方向尺寸**：从 1000.0 改为 20.0

请确认您的选择：
*   如果要应用这些修改，请直接回复“确认”。
*   如果要保留原来的数值，请直接回复“保持原值”。
\item User: 确认
\item Assistant: 配置已完成。当前已确认的项包括：源位置、源能量、源类型以及盒体 z 方向尺寸。
\end{itemize}
\textbf{Internal Chain (Natural Language):}
\begin{itemize}
\item Turn 1: action=summarize\_progress, phase=materials, complete=False.
\item Missing fields: physics.physics\_list, output.format, output.path.
\item LLM used=True; fallback=none; backend=llm\_slot\_frame+runtime\_semantic.
\item Turn 2: action=confirm\_overwrite, phase=source, complete=False.
\item Missing fields: none.
\item LLM used=True; fallback=none; backend=llm\_slot\_frame+runtime\_semantic.
\item Turn 3: action=finalize, phase=physics, complete=True.
\item Missing fields: none.
\item LLM used=True; fallback=none; backend=llm\_semantic\_frame+runtime\_semantic+confirmed\_pending\_overwrite.
\end{itemize}
\subsection{Case MTM08 (education)}
\textbf{Dialogue:}
\begin{itemize}
\item User: For classroom demo, I want gamma attenuation in copper; you can ask only essential missing parts.
\item Assistant: Geometry: still needs geometry structure, box size along x. Materials: updated selected materials; still needs volume-to-material binding. Source: updated particle type; still needs source type, source energy. Physics: updated physics list. Output: still needs output format, output path.
\item User: Okay, copper box 100 mm x 100 mm x 100 mm, point gamma 1 MeV at (0,0,-200) mm, direction +z, physics FTFP\_BERT, output csv.
\item Assistant: This would overwrite an already confirmed value.

Pending change:
Physics list: **Shielding** → **FTFP\_BERT**

Reply **'confirm'** to apply this change, or say **'keep existing'** to keep the current value.
\item User: confirm
\item Assistant: Configuration complete. Confirmed items: physics list, output path, output format, source direction, and source position.
\end{itemize}
\textbf{Internal Chain (Natural Language):}
\begin{itemize}
\item Turn 1: action=summarize\_progress, phase=geometry, complete=False.
\item Missing fields: geometry.structure, geometry.params.module\_x, geometry.params.module\_y, geometry.params.module\_z, materials.volume\_material\_map, source.type, source.energy, source.position, source.direction, output.format, output.path.
\item LLM used=True; fallback=none; backend=llm\_slot\_frame+runtime\_semantic.
\item Turn 2: action=confirm\_overwrite, phase=materials, complete=False.
\item Missing fields: none.
\item LLM used=True; fallback=none; backend=llm\_slot\_frame+runtime\_semantic.
\item Turn 3: action=finalize, phase=source, complete=True.
\item Missing fields: none.
\item LLM used=True; fallback=none; backend=llm\_semantic\_frame+runtime\_semantic+confirmed\_pending\_overwrite.
\end{itemize}
\subsection{Case MTM09 (environmental\_radiation)}
\textbf{Dialogue:}
\begin{itemize}
\item User: 我想先搭一个空气体积中的低能 gamma 传播场景，细节稍后补。
\item Assistant: 目前我们已经完成了几何、材料和源的设置。接下来还需要你补充物理列表、输出格式和输出路径这三项信息。
\item User: 设为 500 mm x 500 mm x 500 mm 的空气盒，材料 G4\_AIR；点源 gamma，能量 0.5 MeV，位置 (0,0,-300) mm，方向 +z；物理列表 QBBC；输出 root。
\item Assistant: 我注意到这会覆盖已经确认的设置。具体修改如下：

- 盒体 x 方向尺寸：1000.0 → 500.0
- 盒体 y 方向尺寸：1000.0 → 500.0

请确认你的选择：
- 如果要应用这些修改，请直接回复“确认”。
- 如果要保留原值，请直接回复“保持原值”。
\item User: 确认
\item Assistant: 配置已完成。当前已确认的设置包括：源位置、源能量、盒体 z 方向尺寸、盒体 y 方向尺寸和盒体 x 方向尺寸。
\end{itemize}
\textbf{Internal Chain (Natural Language):}
\begin{itemize}
\item Turn 1: action=summarize\_progress, phase=materials, complete=False.
\item Missing fields: physics.physics\_list, output.format, output.path.
\item LLM used=True; fallback=none; backend=llm\_slot\_frame+runtime\_semantic.
\item Turn 2: action=confirm\_overwrite, phase=source, complete=False.
\item Missing fields: none.
\item LLM used=True; fallback=none; backend=llm\_slot\_frame+runtime\_semantic.
\item Turn 3: action=finalize, phase=physics, complete=True.
\item Missing fields: none.
\item LLM used=True; fallback=none; backend=llm\_semantic\_frame+runtime\_semantic+confirmed\_pending\_overwrite.
\end{itemize}
\subsection{Case MTM10 (high\_energy\_physics)}
\textbf{Dialogue:}
\begin{itemize}
\item User: Set geometry as a sphere target; I have not decided radius or source details yet.
\item Assistant: You've set the geometry to a sphere. Here's what's complete and what's still needed:

**Complete:**
*   **Geometry:** The structure is a sphere, and you've provided the dimensions for a box that will contain it.

**Still Needed:**
*   **Geometry:** The sphere's radius.
*   **Materials:** Which material to assign to the sphere.
*   **Source:** The type of source, the particle type, and its energy.
*   **Physics:** The physics list to use.
*   **Output:** The format and file path for your results.
\item User: Sphere target radius 80 mm, lead material, electron point source 5 MeV at (0,0,-200) mm toward +z, physics FTFP\_BERT, output hdf5.
\item Assistant: Your configuration is complete. The following items are confirmed: the material is bound to the volume, the output path and format are set, and the physics list is selected.
\end{itemize}
\textbf{Internal Chain (Natural Language):}
\begin{itemize}
\item Turn 1: action=summarize\_progress, phase=geometry, complete=False.
\item Missing fields: geometry.params.child\_rmax, materials.volume\_material\_map, source.type, source.particle, source.energy, source.position, source.direction, physics.physics\_list, output.format, output.path.
\item LLM used=True; fallback=none; backend=llm\_slot\_frame+runtime\_semantic.
\item Turn 2: action=finalize, phase=materials, complete=True.
\item Missing fields: none.
\item LLM used=True; fallback=none; backend=llm\_slot\_frame+runtime\_semantic.
\end{itemize}
\subsection{Case MTM11 (materials\_update)}
\textbf{Dialogue:}
\begin{itemize}
\item User: Create a complete setup: 100 mm x 100 mm x 100 mm copper box, point gamma 1 MeV at (0,0,-200) mm toward +z, physics QBBC, output json.
\item Assistant: Your configuration is complete. The following items have been confirmed: volume-to-material binding, output path, and output format.
\item User: Change the target material to G4\_Al and change the output format to root.
\item Assistant: This change would overwrite two confirmed settings:

*   **Selected materials:** Copper (G4\_Cu) would be replaced with Aluminum (G4\_Al).
*   **Output format:** JSON would be replaced with ROOT.

Please reply **'confirm'** to apply these changes, or **'keep existing'** to keep the current values.
\item User: confirm
\item Assistant: Your configuration is complete. The following items have been confirmed: volume-to-material binding, output path, output format, and selected materials.
\end{itemize}
\textbf{Internal Chain (Natural Language):}
\begin{itemize}
\item Turn 1: action=finalize, phase=materials, complete=True.
\item Missing fields: none.
\item LLM used=True; fallback=none; backend=llm\_slot\_frame+runtime\_semantic.
\item Turn 2: action=confirm\_overwrite, phase=source, complete=False.
\item Missing fields: none.
\item LLM used=True; fallback=none; backend=llm\_slot\_frame+runtime\_semantic.
\item Turn 3: action=finalize, phase=physics, complete=True.
\item Missing fields: none.
\item LLM used=True; fallback=none; backend=llm\_semantic\_frame+runtime\_semantic+confirmed\_pending\_overwrite.
\end{itemize}
\subsection{Case MTM12 (overwrite\_reject)}
\textbf{Dialogue:}
\begin{itemize}
\item User: 请创建一个完整配置：100 mm x 100 mm x 100 mm 的铜盒，材料 G4\_Cu；点源 gamma，能量 1 MeV，位置 (0,0,-200) mm，方向 +z；物理列表 QBBC；输出 json。
\item Assistant: 配置已完成。当前已确认的设置包括：体积与材料绑定、输出路径和输出格式。
\item User: 把材料改成 G4\_Al，并把输出格式改成 root。
\item Assistant: 我注意到这次修改会覆盖已经确认的设置。具体来说：  
- 已选材料将从 `G4\_Cu` 改为 `G4\_Al`  
- 输出格式将从 `json` 改为 `root`  

如果你确定要应用这些修改，请直接回复“确认”；如果想保留原来的值，请直接说“保持原值”。
\item User: 保持原值
\item Assistant: 好的，已保留已选材料（G4\_Cu）和输出格式（json）的原有值，本次覆盖修改没有应用。
\end{itemize}
\textbf{Internal Chain (Natural Language):}
\begin{itemize}
\item Turn 1: action=finalize, phase=materials, complete=True.
\item Missing fields: none.
\item LLM used=True; fallback=none; backend=llm\_slot\_frame+runtime\_semantic.
\item Turn 2: action=confirm\_overwrite, phase=source, complete=False.
\item Missing fields: none.
\item LLM used=True; fallback=none; backend=llm\_slot\_frame+runtime\_semantic.
\item Turn 3: action=reject\_overwrite, phase=physics, complete=True.
\item Missing fields: none.
\item LLM used=True; fallback=none; backend=llm\_semantic\_frame+runtime\_semantic+rejected\_pending\_overwrite.
\end{itemize}
\subsection{Case MTM13 (medical\_detector)}
\textbf{Dialogue:}
\begin{itemize}
\item User: I want a ring of detector modules around a small source region, but I have not fixed the ring radius or module size yet.
\item Assistant: Here's where we stand:

**Geometry:** You've chosen the structure (a torus for the detector ring around a source region). We still need the tube radius of the torus and its major radius.

**Materials:** We need to assign a material to the volume.

**Source:** We still need the source type and the type of particle.

**Physics \& Output:** The physics list, output format, and output path are also still pending.
\item User: Use a ring of 12 modules, each 8 mm x 10 mm x 2 mm, radius 40 mm, clearance 1 mm, material G4\_WATER; point gamma source 0.662 MeV at (0,0,0) mm; physics QBBC; output json.
\item Assistant: This would overwrite an already confirmed value. Pending change: geometry structure: single\_torus -> single\_box; box size along x: None -> 8.0. Reply 'confirm' to apply it, or say 'keep existing' to keep the current value.
\item User: confirm
\item Assistant: Your configuration is complete. The following items have been confirmed: the torus major radius, torus tube radius, volume-to-material binding, and the radius.
\end{itemize}
\textbf{Internal Chain (Natural Language):}
\begin{itemize}
\item Turn 1: action=summarize\_progress, phase=geometry, complete=False.
\item Missing fields: geometry.params.torus\_rmax, geometry.params.torus\_rtor, materials.volume\_material\_map, source.type, source.particle, source.energy, source.position, source.direction, physics.physics\_list, output.format, output.path.
\item LLM used=True; fallback=none; backend=llm\_slot\_frame+runtime\_semantic.
\item Turn 2: action=confirm\_overwrite, phase=geometry, complete=False.
\item Missing fields: geometry.params.torus\_rmax, geometry.params.torus\_rtor.
\item LLM used=True; fallback=none; backend=llm\_slot\_frame+runtime\_semantic.
\item Turn 3: action=finalize, phase=materials, complete=True.
\item Missing fields: none.
\item LLM used=True; fallback=none; backend=llm\_semantic\_frame+runtime\_semantic+confirmed\_pending\_overwrite.
\end{itemize}
\subsection{Case MTM14 (industrial\_panel)}
\textbf{Dialogue:}
\begin{itemize}
\item User: 我想做一个二维阵列探测板，但 pitch 和材料还没定。
\item Assistant: 目前我们已经确定了使用二维阵列探测板作为几何结构，并完成了盒体尺寸的设置。

接下来还需要补充以下信息：
- **几何**：阵列模块的尺寸、x 和 y 方向的阵列数量。
- **材料**：每个体积对应的材料。
- **源**：源的类型、粒子种类和能量。
- **物理**：需要使用的物理列表。
- **输出**：输出文件的格式和保存路径。

请随时告诉我这些信息，我们可以一项项来确认。
\item User: 设成 4 x 3 的盒体阵列，每个模块 10 mm x 10 mm x 2 mm，pitch\_x 12 mm，pitch\_y 14 mm，间隙 1 mm，材料 G4\_Si；点源 gamma 0.511 MeV，位置 (0,0,-80) mm，方向 +z；物理列表 QBBC；输出 root。
\item Assistant: 我检测到这会覆盖已经确认的值。待修改内容：几何结构类型：grid -> single\_box；geometry.chosen\_skeleton：grid\_modules -> single\_box。如果要应用这些修改，请直接回复“确认”；如果要保留原值，请直接说“保持原值”。
\item User: 确认
\item Assistant: 配置已完成。当前已确认的几何参数包括：几何结构类型、盒体在 x、y、z 三个方向的尺寸。
\end{itemize}
\textbf{Internal Chain (Natural Language):}
\begin{itemize}
\item Turn 1: action=summarize\_progress, phase=geometry, complete=False.
\item Missing fields: geometry.graph\_program, materials.volume\_material\_map, source.type, source.particle, source.energy, source.position, source.direction, physics.physics\_list, output.format, output.path, geometry.ask.grid.module\_size, geometry.ask.grid.count\_x, geometry.ask.grid.count\_y, geometry.ask.grid.pitch\_x, geometry.ask.grid.pitch\_y, geometry.ask.grid.clearance.
\item LLM used=True; fallback=none; backend=llm\_semantic\_frame+runtime\_semantic.
\item Turn 2: action=confirm\_overwrite, phase=materials, complete=False.
\item Missing fields: none.
\item LLM used=True; fallback=none; backend=llm\_slot\_frame+runtime\_semantic.
\item Turn 3: action=finalize, phase=source, complete=True.
\item Missing fields: none.
\item LLM used=True; fallback=none; backend=llm\_semantic\_frame+runtime\_semantic+confirmed\_pending\_overwrite.
\end{itemize}
\subsection{Case MTM15 (shielding\_stack)}
\textbf{Dialogue:}
\begin{itemize}
\item User: Build a stacked absorber sandwich along z; I will provide the layer sizes next.
\item Assistant: You've set up the geometry structure for a stacked absorber sandwich. Here's what's still needed to complete the setup:

**Geometry:** We need the inner radii (r1, r2, r3) and the layer positions (z1, z2) for the sandwich.
**Materials:** We need to assign materials to the volumes.
**Source:** We need the source type, particle type, and energy.
**Physics:** We need to select a physics list.
**Output:** We need the output format and file path.
\item User: Use stacked layers with footprint 20 mm x 20 mm, thicknesses 2 mm, 4 mm, 6 mm, layer clearance 0.5 mm, outer box 30 mm x 30 mm x 20 mm, container clearance 1 mm, material G4\_Al; beam gamma 1.25 MeV from (0,0,-120) mm along +z; physics Shielding; output csv.
\item Assistant: Here's your current setup:

**Geometry:** Box size is set. We still need the inner radii (r1, r2, r3) and the layer positions (z1, z2) for the sandwich layers.
**Materials:** Materials are selected and assigned.
**Source:** Source type, particle, and energy are set.
**Physics:** Physics list is selected.
**Output:** Output format and file path are set.
\end{itemize}
\textbf{Internal Chain (Natural Language):}
\begin{itemize}
\item Turn 1: action=summarize\_progress, phase=geometry, complete=False.
\item Missing fields: geometry.params.r1, geometry.params.r2, geometry.params.r3, geometry.params.z1, geometry.params.z2, geometry.params.z3, materials.volume\_material\_map, source.type, source.particle, source.energy, source.position, source.direction, physics.physics\_list, output.format, output.path.
\item LLM used=True; fallback=none; backend=llm\_slot\_frame+runtime\_semantic.
\item Turn 2: action=summarize\_progress, phase=geometry, complete=False.
\item Missing fields: geometry.params.r1, geometry.params.r2, geometry.params.r3, geometry.params.z1, geometry.params.z2, geometry.params.z3.
\item LLM used=True; fallback=none; backend=llm\_slot\_frame+runtime\_semantic.
\end{itemize}
\subsection{Case MTM16 (nested\_target)}
\textbf{Dialogue:}
\begin{itemize}
\item User: 我想在一个盒子里嵌一个圆柱靶，先别急着补全。
\item Assistant: 目前我们已经确定了使用盒体作为几何结构，并设置了盒体的尺寸。接下来还需要补充圆柱的半长和半径、为体积指定材料、选择源类型和粒子类型，以及配置物理列表和输出选项。
\item User: 外盒 120 mm x 120 mm x 120 mm，内嵌铜圆柱半径 20 mm、半长 40 mm，嵌套间隙 2 mm；材料 G4\_Cu；束流 gamma 1 MeV，位置 (0,0,-150) mm，方向 +z；物理列表 QBBC；输出 json。
\item Assistant: 我检测到这会覆盖已经确认的值。待修改内容：几何结构类型：single\_tubs -> single\_box；盒体 x 方向尺寸：None -> 120.0。如果要应用这些修改，请直接回复“确认”；如果要保留原值，请直接说“保持原值”。
\item User: 确认
\item Assistant: 好的，我们已同步了您设置的几何结构类型和盒体尺寸。

接下来，为了完成配置，我们还需要：
1.  为体积指定材料。
2.  设置源的参数（类型、粒子、能量和位置）。
3.  选择物理列表。
4.  确定输出格式和路径。
\end{itemize}
\textbf{Internal Chain (Natural Language):}
\begin{itemize}
\item Turn 1: action=summarize\_progress, phase=geometry, complete=False.
\item Missing fields: geometry.params.child\_hz, geometry.params.child\_rmax, materials.volume\_material\_map, source.type, source.particle, source.energy, source.position, source.direction, physics.physics\_list, output.format, output.path.
\item LLM used=True; fallback=none; backend=llm\_slot\_frame+runtime\_semantic.
\item Turn 2: action=confirm\_overwrite, phase=geometry, complete=False.
\item Missing fields: geometry.params.child\_hz, geometry.params.child\_rmax, materials.volume\_material\_map, source.type, source.particle, source.energy, source.position, source.direction, physics.physics\_list, output.format, output.path.
\item LLM used=True; fallback=none; backend=llm\_slot\_frame+runtime\_semantic.
\item Turn 3: action=summarize\_progress, phase=materials, complete=False.
\item Missing fields: materials.volume\_material\_map, source.type, source.particle, source.energy, source.position, source.direction, physics.physics\_list, output.format, output.path.
\item LLM used=True; fallback=none; backend=llm\_semantic\_frame+runtime\_semantic+confirmed\_pending\_overwrite.
\end{itemize}
\subsection{Case MTM17 (coaxial\_shield)}
\textbf{Dialogue:}
\begin{itemize}
\item User: I need a shell-style shielding toy model, but the shell thicknesses are still open.
\item Assistant: Here's a summary of your setup so far:

**Completed:**
- **Geometry:** Structure set to a shell-style shielding model. Box dimensions are confirmed.
- **Physics:** Physics list is confirmed.

**Still Needed:**
- **Geometry:** The shell's radius.
- **Materials:** Assignment of materials to each volume.
- **Source:** Type, particle type, and energy.
- **Output:** Format and file path.
\item User: Use a shell with inner radius 15 mm, thicknesses 5 mm and 8 mm, half length 40 mm, nested core radius 10 mm, nested core half length 30 mm, clearance 1 mm, material G4\_Pb; isotropic neutron source 2 MeV at (0,0,0) mm; physics Shielding; output hdf5.
\item Assistant: This would overwrite an already confirmed value. Pending change: geometry structure: single\_sphere -> single\_polycone; radius: None -> 15.0. Reply 'confirm' to apply it, or say 'keep existing' to keep the current value.
\item User: confirm
\item Assistant: Here's a summary of your setup so far:

**Completed:**
- **Geometry:** Updated the structure, radius, and half length. The radii (r1, r2, r3) are confirmed.
- **Materials:** Updated the assignment of materials to volumes.
- **Source:** Confirmed the source direction and position.
- **Output:** Confirmed the file path and format.

**Still Needed:**
- **Geometry:** The z-positions (z1, z2, z3) for the shells.
\end{itemize}
\textbf{Internal Chain (Natural Language):}
\begin{itemize}
\item Turn 1: action=summarize\_progress, phase=geometry, complete=False.
\item Missing fields: geometry.params.child\_rmax, materials.volume\_material\_map, source.type, source.particle, source.energy, source.position, source.direction, output.format, output.path.
\item LLM used=True; fallback=none; backend=llm\_slot\_frame+runtime\_semantic.
\item Turn 2: action=confirm\_overwrite, phase=geometry, complete=False.
\item Missing fields: geometry.params.child\_rmax.
\item LLM used=True; fallback=none; backend=llm\_slot\_frame+runtime\_semantic.
\item Turn 3: action=summarize\_progress, phase=geometry, complete=False.
\item Missing fields: geometry.params.z1, geometry.params.z2, geometry.params.z3.
\item LLM used=True; fallback=none; backend=llm\_semantic\_frame+runtime\_semantic+confirmed\_pending\_overwrite.
\end{itemize}
\subsection{Case MTM18 (boolean\_cavity)}
\textbf{Dialogue:}
\begin{itemize}
\item User: 我想做一个挖空的盒体几何，但先只给你意图。
\item Assistant: 目前我们已确定了使用挖空的盒体作为几何结构，接下来还需要补充一些具体信息：

- **几何**：需要盒体在 x、y、z 三个方向上的尺寸。
- **材料**：需要指定体积与材料的绑定关系。
- **源**：需要选择源的类型（例如粒子类型、能量等）。

请随时告诉我这些信息，我们可以继续推进配置。
\item User: 做减法布尔体：A 为 40 mm x 40 mm x 40 mm 盒体，B 为 20 mm x 20 mm x 20 mm 盒体，从 A 中减去 B；材料 G4\_Al；点源 gamma 0.3 MeV，位置 (0,0,-60) mm，方向 +z；物理列表 QBBC；输出 xml。
\item Assistant: 配置已完成。当前已确认的设置包括：体积与材料绑定、输出路径、输出格式和物理列表。
\end{itemize}
\textbf{Internal Chain (Natural Language):}
\begin{itemize}
\item Turn 1: action=summarize\_progress, phase=geometry, complete=False.
\item Missing fields: geometry.params.module\_x, geometry.params.module\_y, geometry.params.module\_z, materials.volume\_material\_map, source.type, source.particle, source.energy, source.position, source.direction, physics.physics\_list, output.format, output.path.
\item LLM used=True; fallback=none; backend=llm\_slot\_frame+runtime\_semantic.
\item Turn 2: action=finalize, phase=materials, complete=True.
\item Missing fields: none.
\item LLM used=True; fallback=none; backend=llm\_slot\_frame+runtime\_semantic.
\end{itemize}
\subsection{Case MTM19 (ring\_complete)}
\textbf{Dialogue:}
\begin{itemize}
\item User: Create a complete ring setup: 16 modules, each 6 mm x 6 mm x 2 mm, radius 32 mm, clearance 0.5 mm, material G4\_Si; point gamma source 0.511 MeV at (0,0,0) mm; physics QBBC; output json.
\item Assistant: Here's your current setup status:

**Geometry:** I've set the structure and the module's X and Y dimensions. I still need the ring's overall radius and the clearance between modules.

**Materials, Source, Physics, and Output:** All have been updated with the details you provided.
\end{itemize}
\textbf{Internal Chain (Natural Language):}
\begin{itemize}
\item Turn 1: action=summarize\_progress, phase=materials, complete=False.
\item Missing fields: geometry.ask.ring.radius, geometry.ask.ring.clearance.
\item LLM used=True; fallback=none; backend=llm\_slot\_frame+runtime\_semantic.
\end{itemize}
\subsection{Case MTM20 (grid\_complete)}
\textbf{Dialogue:}
\begin{itemize}
\item User: 请创建完整的阵列配置：3 x 3 盒体阵列，每个模块 12 mm x 12 mm x 3 mm，pitch\_x 15 mm，pitch\_y 15 mm，间隙 1 mm，材料 G4\_CsI；点源 gamma 0.662 MeV，位置 (0,0,-90) mm，方向 +z；物理列表 QBBC；输出 root。
\item Assistant: 配置已完成！当前已确认的设置包括：体积与材料绑定、输出路径、y方向pitch和x方向pitch。
\end{itemize}
\textbf{Internal Chain (Natural Language):}
\begin{itemize}
\item Turn 1: action=finalize, phase=materials, complete=True.
\item Missing fields: none.
\item LLM used=True; fallback=none; backend=llm\_slot\_frame+runtime\_semantic.
\end{itemize}
\subsection{Case MTM21 (layered\_detector)}
\textbf{Dialogue:}
\begin{itemize}
\item User: I want a stacked detector slab, but the individual layer thicknesses are not fixed yet.
\item Assistant: Geometry: updated geometry structure; still needs stack layer footprint, stack layer thicknesses. Materials: still needs volume-to-material binding. Source: still needs source type, particle type. Physics: still needs physics list. Output: still needs output format, output path.
\item User: Use stack footprint 25 mm x 25 mm, thicknesses 1 mm, 2 mm, 1 mm, layer clearance 0.2 mm, outer box 30 mm x 30 mm x 12 mm, container clearance 0.5 mm, material G4\_Si; beam electron 5 MeV from (0,0,-50) mm along +z; physics FTFP\_BERT; output json.
\item Assistant: This would overwrite an already confirmed value. Pending change: geometry structure: stack -> single\_box; geometry.chosen\_skeleton: stack\_in\_box -> single\_box. Reply 'confirm' to apply it, or say 'keep existing' to keep the current value.
\item User: confirm
\item Assistant: Configuration complete. Confirmed items: box size along z, box size along y, box size along x, and geometry structure.
\end{itemize}
\textbf{Internal Chain (Natural Language):}
\begin{itemize}
\item Turn 1: action=summarize\_progress, phase=geometry, complete=False.
\item Missing fields: geometry.graph\_program, materials.volume\_material\_map, source.type, source.particle, source.energy, source.position, source.direction, physics.physics\_list, output.format, output.path, geometry.ask.stack.footprint, geometry.ask.stack.thicknesses, geometry.ask.stack.layer\_clearance, geometry.ask.stack.parent\_size, geometry.ask.stack.nest\_clearance.
\item LLM used=True; fallback=none; backend=llm\_semantic\_frame+runtime\_semantic.
\item Turn 2: action=confirm\_overwrite, phase=materials, complete=False.
\item Missing fields: none.
\item LLM used=True; fallback=none; backend=llm\_slot\_frame+runtime\_semantic.
\item Turn 3: action=finalize, phase=source, complete=True.
\item Missing fields: none.
\item LLM used=True; fallback=none; backend=llm\_slot\_frame+runtime\_semantic+confirmed\_pending\_overwrite.
\end{itemize}
\subsection{Case MTM22 (embedded\_sensor)}
\textbf{Dialogue:}
\begin{itemize}
\item User: 请创建完整的嵌套配置：外盒 80 mm x 80 mm x 80 mm，内嵌铅圆柱半径 15 mm、半长 25 mm，嵌套间隙 1 mm；材料 G4\_Pb；点源 gamma 1 MeV，位置 (0,0,-100) mm，方向 +z；物理列表 Shielding；输出 csv。
\item Assistant: 你提到的外盒是长方体，那它的另外两个方向的尺寸分别是多少呢？
\end{itemize}
\textbf{Internal Chain (Natural Language):}
\begin{itemize}
\item Turn 1: action=ask\_clarification, phase=geometry, complete=False.
\item Missing fields: geometry.structure, geometry.params.module\_x, geometry.params.module\_y, geometry.params.module\_z, materials.volume\_material\_map, source.type, source.particle, source.energy, source.position, source.direction, physics.physics\_list, output.format, output.path.
\item LLM used=True; fallback=none; backend=llm\_slot\_frame+runtime\_semantic.
\end{itemize}
\subsection{Case MTM23 (layered\_shield\_complete)}
\textbf{Dialogue:}
\begin{itemize}
\item User: Create a complete shell configuration: inner radius 12 mm, thicknesses 4 mm and 6 mm, half length 35 mm, nested core radius 8 mm, nested core half length 25 mm, clearance 1 mm, material G4\_Fe; isotropic gamma source 0.8 MeV at (0,0,0) mm; physics QBBC; output hdf5.
\item Assistant: Configuration complete. All geometry, material, source, physics, and output settings are now finalized.
\end{itemize}
\textbf{Internal Chain (Natural Language):}
\begin{itemize}
\item Turn 1: action=finalize, phase=materials, complete=True.
\item Missing fields: none.
\item LLM used=True; fallback=none; backend=llm\_slot\_frame+runtime\_semantic.
\end{itemize}
\subsection{Case MTM24 (boolean\_complete)}
\textbf{Dialogue:}
\begin{itemize}
\item User: 请创建完整布尔配置：做并集布尔体，A 为 30 mm x 30 mm x 20 mm 盒体，B 为 10 mm x 10 mm x 40 mm 盒体，与 A 做 union；材料 G4\_Al；束流 gamma 0.5 MeV，位置 (0,0,-70) mm，方向 +z；物理列表 QBBC；输出 json。
\item Assistant: 配置已完成！当前已确认的设置包括：体积与材料绑定、输出路径，以及布尔体 B 的 x、y、z 尺寸。
\end{itemize}
\textbf{Internal Chain (Natural Language):}
\begin{itemize}
\item Turn 1: action=finalize, phase=materials, complete=True.
\item Missing fields: none.
\item LLM used=True; fallback=none; backend=llm\_slot\_frame+runtime\_semantic.
\end{itemize}
\end{document}